\documentclass[11pt,a4paper]{article}

\usepackage[margin=1in]{geometry}
\usepackage[T1]{fontenc}
\usepackage[utf8]{inputenc}
\usepackage{lmodern}
\usepackage{microtype}
\usepackage{amsmath,amssymb,amsthm,mathtools}
\usepackage{enumitem}
\usepackage{booktabs}
\usepackage{longtable}
\usepackage{xcolor}
\usepackage{hyperref}
\usepackage{tikz}
\usetikzlibrary{arrows.meta,positioning,fit,calc}

\hypersetup{colorlinks=true,linkcolor=blue!50!black,urlcolor=blue!50!black,citecolor=blue!50!black}
\setlist[itemize]{leftmargin=1.5em,itemsep=0.2em,topsep=0.3em}
\setlist[enumerate]{leftmargin=1.7em,itemsep=0.2em,topsep=0.3em}

\newtheorem{definition}{Definition}[section]
\newtheorem{theorem}[definition]{Theorem}
\newtheorem{lemma}[definition]{Lemma}
\newtheorem{proposition}[definition]{Proposition}
\newtheorem{corollary}[definition]{Corollary}
\newtheorem{obligation}[definition]{Obligation}
\theoremstyle{remark}
\newtheorem{remark}[definition]{Remark}
\newtheorem{firewall}[definition]{Firewall}

\newcommand{\T}{\mathcal T}
\newcommand{\K}{\mathcal K}
\newcommand{\Hone}{\mathcal H_1}
\newcommand{\Hz}{\mathcal H_Z}
\newcommand{\Cfg}{\mathcal C}
\newcommand{\Val}{\mathsf{Val}}
\newcommand{\MDL}{\operatorname{MDL}}
\newcommand{\Argmin}{\operatorname*{Argmin}}
\newcommand{\PT}{\operatorname{PT}}
\newcommand{\Atom}{\mathsf{Atom}}
\newcommand{\Elem}{\mathsf{Element}}
\newcommand{\Nat}{\mathbb N}
\newcommand{\Rphys}{\mathcal R_{\rm phys}}
\newcommand{\qchem}{q_{\rm chem}}
\newcommand{\simchem}{\sim_{\rm chem}}
\newcommand{\simper}{\sim_{\rm per}}
\newcommand{\defeq}{\vcentcolon=}

\title{DASHI Atomic and Periodic-Table Formalism\\
\large Kernel Filling, MDL Selection, Valence Recurrence, and Provenance Gates}
\author{Prepared from the attached \emph{DASHI Atom} archive and current DASHI formalism notes}
\date{2026-09-11}

\begin{document}
\maketitle

\begin{abstract}
This note writes the atom/chemistry formalism as a typed generative theory.  The
core claim is not that a finite geometric chart can be counted and renamed as
chemistry.  The core claim is that a ternary/369 structural substrate can feed a
realization layer; the realization layer supplies a discrete fermionic kernel;
hard exclusion and a minimum-description-length ordering select stable fillings;
closed kernels yield noble-gas-like closures; and periodic-table structure is the
quotient of the atomic family by recurrent boundary/valence classes.  Provenance
is kept append-only: archive statements, Python-test claims, Agda owners, and
external comparison models have different authority.
\end{abstract}

\tableofcontents
\newpage

\section{Source register and provenance discipline}

The attached archive contains several distinct source roles.  The formalism below
keeps them separated, because using the same mathematics under the wrong source
role is precisely the type error this lane is designed to avoid.

\begin{longtable}{p{0.22\linewidth}p{0.37\linewidth}p{0.31\linewidth}}
\toprule
\textbf{Source layer} & \textbf{Content imported} & \textbf{Authority granted here}\\
\midrule
\endhead
Attached \emph{DASHI Atom} conversation archive & Atoms-as-dictionaries, noble gases as saturated kernels, MDL/exclusion filling, and the first closures $Z=2,10,18$ in the remembered Python-test lane. & Historical/provenance evidence for the internal DASHI theory and its conversational/test lineage.\\
Attached rendered pages 294--301 & The atom-centric reconstruction: kernel states $(n,\ell,m,s)$, exclusion quotient, MDL proxy $n+\alpha\ell$, closures, boundary-condition equivalence, magic-number analogy, and explicit non-claims. & Formal content to be stabilized into the typed theory below.\\
Attached page 418 & Move from topographic closure tests to a Hamiltonian/ionization-energy validation path, with $IE_1(Z)=E(Z-1)-E(Z)$ and scale as derivative. & Snowball target for the next empirical/scale bridge, not a completed physical calibration theorem.\\
Current markdown formalism & Clean typed presentation: ternary carrier, atomic species, Hilbert/exterior-power layer, configuration minimization, valence recurrence, periodic quotient, and non-collapse laws. & Base text for the rendered LaTeX formalism.\\
Agda/repo owners & Names such as \texttt{AtomicFermionShellExact}, \texttt{AtomicGenerationPipelineExact}, \texttt{AtomicShellRecurrence}, and \texttt{AtomicPeriodicTableRecoveryBoundary}. & Same-object anchors and implementation targets; they do not by themselves promote unresolved physical gates.\\
External topological periodic-matrix claims & $S^4\times T^3$, $86+86$, $Z=172$, $1/C(n)$, and arithmetic/discriminant claims. & Candidate external chart only.  No same-object identity with DASHI is assumed.\\
\bottomrule
\end{longtable}

\begin{firewall}[Append-only provenance]
Acquisition order may be out of dependency order, but payment order may not be.
Later evidence may be retained as a snowball target, but it does not rewrite the
historical source state and does not pay an earlier identity, scale, or empirical
adequacy obligation.
\end{firewall}

\section{Executive claim}

The DASHI atom formalism is the following compositional statement:
\[
\boxed{\begin{gathered}
\text{369/ternary structure}\to\text{realized fermionic kernel}\to\text{MDL-ordered exclusion filling}\\
\to\text{closed-shell recurrence}\to\text{periodic-table quotient}
\end{gathered}}
\]

This is stronger and cleaner than a finite-table embedding.  It generates a
parameterized family indexed by nuclear charge and then quotients that family by
chemistry-visible recurrence.  It does not count cells in a manifold and assign
names afterward.

\section{Ternary substrate}

\begin{definition}[Ternary carrier]
For each finite depth $n$, let
\[
\T_n = \{-1,0,+1\}^n
\]
or equivalently $\{0,1,2\}^n$ after relabelling.  Let $\iota:\T_n\to\T_n$ be
an involution, written $\iota(x)=-x$ in the balanced representation, with fixed
neutral symbol $0$.
\end{definition}

\begin{definition}[Agreement ultrametric]
Let $\ell(x,y)$ be the greatest common prefix/agreement depth of $x,y\in\T_n$.
For fixed $0<\rho<1$, define
\[
 d_{\T}(x,y)=\rho^{\ell(x,y)}.
\]
This metric records structural agreement depth, not physical distance.
\end{definition}

\begin{definition}[Admissible contraction]
An admissible structural evolution is a map $K:\T_n\to\T_n$ together with a
certificate
\[
 d_{\T}(Kx,Ky)\le \lambda d_{\T}(x,y),\qquad 0\le \lambda<1.
\]
A stable structural class is either a fixed point $Kx^\star=x^\star$ or a
quotient-invariant orbit $[Kx]_{\sim}=[x]_{\sim}$.
\end{definition}

\begin{firewall}[No numerological promotion]
\[
 |\T_n|=N \quad\not\Rightarrow\quad N\text{ chemical elements}.
\]
The ternary/369 carrier is a structural substrate.  Atomic meaning requires a
realization map and the downstream fermionic and energetic witnesses.
\end{firewall}

\section{Physical realization and kernel states}

\begin{definition}[Realization map]
A physical realization is a typed map
\[
 \Phi:\T_n\longrightarrow \Rphys
\]
that preserves the structural data required by its consumer.  Typical required
laws are equivariance
\[
 \Phi(Kx)=U\Phi(x)
\]
for a physical evolution $U$, and quotient preservation
\[
 x\sim_{\T}y \Longrightarrow \Phi(x)\sim_{\rm phys}\Phi(y).
\]
\end{definition}

\begin{definition}[Kernel mode]
The realized atomic kernel contains fermionic modes labelled by
\[
 k=(n,\ell,m,s),
\]
where
\[
 n\in\Nat_{>0},\qquad \ell\in\{0,\ldots,n-1\},\qquad
 m\in\{-\ell,\ldots,+\ell\},\qquad s\in\{\uparrow,\downarrow\}.
\]
The one-particle carrier is
\[
 \Hone = \bigoplus_{n,\ell} V_\ell\otimes V_{1/2}.
\]
\end{definition}

\begin{proposition}[Subshell and shell capacities]
For fixed angular index $\ell$,
\[
 g(n,\ell)=2(2\ell+1).
\]
Hence
\[
 g_s=2,\qquad g_p=6,
\qquad g_d=10,
\qquad g_f=14.
\]
For a principal shell $n$,
\[
 G_n = \sum_{\ell=0}^{n-1}2(2\ell+1)=2n^2.
\]
Thus
\[
 G_1=2,
\quad G_2=8,
\quad G_3=18,
\quad G_4=32,
\quad \ldots
\]
\end{proposition}

\begin{firewall}[Representation before capacity]
The capacity theorem depends on a supplied rotational/spin representation.
Therefore
\[
\text{triadic carrier}\not\Rightarrow(n,\ell,m,s)
\]
without a realization/representation bridge.
\end{firewall}

\section{Atomic identity}

\begin{definition}[Atomic species]
An atomic species is a triple
\[
 A=(Z,N,N_e)\in\Nat^3,
\]
where $Z$ is proton number, $N$ is neutron number, and $N_e$ is electron number.
Neutrality is
\[
 \mathsf{Neutral}(A)\;\Longleftrightarrow\;N_e=Z.
\]
Element identity is indexed by nuclear charge:
\[
 A_1\sim_{\rm element}A_2 \Longleftrightarrow Z(A_1)=Z(A_2).
\]
Isotope identity additionally fixes neutron number:
\[
 A_1\sim_{\rm isotope}A_2 \Longleftrightarrow Z_1=Z_2\land N_1=N_2.
\]
\end{definition}

\begin{firewall}[Electronic state does not define element]
Two atoms with the same $Z$ may have different electronic configurations.  An
electronic filling pattern is therefore not the primitive identity of the
element.
\end{firewall}

\section{Hard exclusion and the many-fermion carrier}

\begin{definition}[Exclusion quotient]
The Pauli/exclusion rule is treated as a hard admissibility quotient on the
kernel:
\[
 n_k\in\{0,1\}\qquad\text{for every mode }k=(n,\ell,m,s).
\]
Equivalently, no two indistinguishable fermions may occupy the same key in the
same kernel scope.
\end{definition}

\begin{definition}[Many-electron carrier]
For electron number $N_e$, the electronic state lies in the antisymmetric sector
\[
 \mathcal H_{N_e}=\bigwedge^{N_e}\Hone,
\]
not in the raw tensor product $\Hone^{\otimes N_e}$.
\end{definition}

\begin{corollary}[Finite occupancy]
For any subshell,
\[
 0\le N_{n\ell}\le 2(2\ell+1).
\]
A subshell is closed when
\[
 N_{n\ell}=2(2\ell+1),
\]
and a principal shell is closed when
\[
 N_n=2n^2.
\]
\end{corollary}

\section{MDL filling}

\begin{definition}[MDL cost]
The simplest programmatic filling cost used in the archive is
\[
 C_{\alpha}(n,\ell)=n+\alpha\ell,
 \qquad 0<\alpha<1.
\]
This is a description-cost proxy.  It is not asserted to be the physical
Hamiltonian energy.
\end{definition}

\begin{definition}[Configuration]
A finite electronic configuration is an occupancy function
\[
 C:k\longmapsto n_k\in\{0,1\}
\]
satisfying
\[
 \sum_k n_k=N_e.
\]
The MDL-selected configuration is any minimizer
\[
 C^{\rm MDL}_{Z,N_e}\in
 \Argmin_{C\in\Cfg(Z,N_e)} \MDL(C),
\]
where, in the simplest block model, $\MDL(C)$ is generated from
$C_{\alpha}(n,\ell)$.
\end{definition}

\begin{proposition}[Finite filling algorithm]
At fixed cutoff:
\begin{enumerate}
\item enumerate all blocks $(n,\ell)$;
\item assign cost $C_{\alpha}(n,\ell)$;
\item sort by that cost;
\item expose exactly $g(n,\ell)$ fermionic keys in each block;
\item fill admissible keys until cumulative occupancy is $N_e$.
\end{enumerate}
Thus the structural simulation is
\[
\boxed{\text{MDL ordering}+\text{exclusion}+\text{degeneracy}.}
\]
\end{proposition}

\begin{definition}[Closure]
A filling $C$ is closed at MDL depth $d$ when all admissible kernel modes of
cost at most $d$ are occupied and any admissible extension requires a strictly
larger cost class:
\[
 \forall e\in\operatorname{AdmExt}(C),
 \qquad \MDL(C\oplus e)>\MDL(C).
\]
\end{definition}

\begin{remark}[Historical finite regression]
The archive records the first structural closures as
\[
 Z=2,\qquad Z=10,\qquad Z=18,
\]
corresponding respectively to the saturated $1s$, $2s+2p$, and $3s+3p$
regimes in the simplified filling model.  These values are retained as
historical regression targets, not elevated into an independent proof of all
physical atomic spectra.
\end{remark}

\begin{firewall}[Cost is not spectrum]
\[
 C_{\alpha}(n,\ell)\not\equiv E^{\rm physical}_{n\ell}.
\]
In particular, the simplified MDL rule can expose ordering sensitivity near
competing subshells.  A physical spectrum requires the Hamiltonian/scale bridge
below.
\end{firewall}

\section{Atoms as dictionaries}

\begin{definition}[Atomic dictionary]
An atom is represented as a typed dictionary
\[
 A=\{\K_A,B_A,G_A,L_A\},
\]
where
\begin{itemize}
\item $\K_A$ is the internal fermionic kernel;
\item $B_A$ is the writable/exposed interaction boundary;
\item $G_A$ is the stabilizer/symmetry data;
\item $L_A$ is the description/selection cost.
\end{itemize}
The dictionary language is a quotient/interface description of the same
occupancy state, not a replacement for the underlying fermionic carrier.
\end{definition}

\begin{definition}[Saturated dictionary]
A dictionary is saturated when its active kernel classes are closed and no
low-cost writable boundary remains.  Equivalently, every admissible extension
raises the relevant description cost.
\end{definition}

\begin{proposition}[Noble-gas structural class]
A noble-gas-like structural state is a saturated dictionary whose valence
boundary is closed.  In the archive's terminology this gives maximal local
symmetry, minimal writable interaction boundary, and MDL resistance to
extension.
\end{proposition}

This is the precise mathematical content of the ``perfectly packed suitcase''
analogy: saturation is not a visual metaphor but a certificate that no
lower-description admissible extension remains in the current kernel class.

\section{Valence and periodicity}

\begin{definition}[Active boundary]
Given a physical or proxy ordering, let $A_C(k)$ state that mode $k$ is active
at the interaction boundary of configuration $C$.  The boundary projection is
\[
 V(C)=\{k\mid A_C(k)\}.
\]
For an actual Hamiltonian one may define activity through an energy window
$W_{\Delta}(E_F)$; for the structural MDL model it is the currently writable
minimal-cost boundary.
\end{definition}

\begin{definition}[Valence class]
Let $\simchem$ identify boundary states with the same chemistry-visible
structure.  Define
\[
 \Val(C)=[V(C)]_{\simchem}.
\]
\end{definition}

\begin{definition}[Shell-completion step]
A relation
\[
 C\leadsto C'
\]
means that $C'$ is reached from $C$ through an admissible shell-completion
progression while preserving the selected structural laws.
\end{definition}

\begin{definition}[Periodicity]
Periodicity is recurrence of valence classes across completion events:
\[
 C\leadsto C'\land\operatorname{Closed}(C)
 \Longrightarrow
 R_{\rm val}(\Val(C),\Val(C')).
\]
It is not postulated as a single fixed translation $V(Z+p)=V(Z)$ for every
$Z$.
\end{definition}

\begin{definition}[Periodic-table quotient]
Let
\[
 \mathcal P=\{\Atom(Z):Z\in\Nat_{>0}\}
\]
be the generated neutral atomic family.  For a periodic observation map
\[
 Q_{\rm per}:\mathcal P\to\mathcal O_{\rm per},
\]
define
\[
 A_Z\simper A_{Z'}
 \Longleftrightarrow
 Q_{\rm per}(A_Z)=Q_{\rm per}(A_{Z'}).
\]
The periodic table is the ordered quotient
\[
 \boxed{\PT=\mathcal P/\simper}
\]
together with nuclear-charge order $Z_1<Z_2$.
\end{definition}

\begin{remark}[Groups as boundary equivalence]
The archive's example $Z=3$ and $Z=11$ illustrates the intended notion: both
have one active electron outside a saturated inner kernel in the simplified
model.  The repeated chemistry class is therefore a boundary-equivalence class,
not an empirical coincidence inserted as a lookup table.
\end{remark}

\section{Molecules as dictionary merges}

A molecule is not required as a new primitive.  Given atomic dictionaries
$A_1,\dots,A_r$, define a candidate merge
\[
 M=A_1\oplus\cdots\oplus A_r\oplus S,
\]
where $S$ is a set of shared boundary keys.  A structural bond is admitted when
\[
 \operatorname{Adm}(M)
\]
and, in the MDL model,
\[
 \MDL(M)<\sum_i\MDL(A_i).
\]
This is a clean place to cross-pollinate the archive's proposed interpretation
of covalent/ionic alternatives as different merge/quotient choices.  It remains
a chemistry-model interface until a physical bonding-energy interpretation is
supplied.

\section{Hamiltonian and scale bridge}

The MDL constructor produces structural ordering.  A physical atom requires a
Hamiltonian realization.  For neutral charge $Z$ define schematically
\[
 H_Z=\sum_{i=1}^{Z}h_Z(i)+\sum_{i<j}V_{ee}(i,j),
\]
with the familiar Coulomb realization only when separately adopted:
\[
 h_Z=-\frac{\hbar^2}{2m_e}\nabla^2-
 \frac{Ze^2}{4\pi\varepsilon_0 r},
 \qquad
 V_{ee}(i,j)=\frac{e^2}{4\pi\varepsilon_0|r_i-r_j|}.
\]

The physical ground state is
\[
 \Psi_Z^\star\in
 \Argmin_{\Psi\in\bigwedge^Z\Hone}
 \langle\Psi,H_Z\Psi\rangle.
\]
The structural and physical orderings must then be related by an explicit
bridge, rather than silently identified.

\begin{definition}[Ionization observable]
Once a calibrated total-energy function $E(Z,N_e)$ is supplied,
\[
 IE_1(Z)=E(Z,Z-1)-E(Z,Z).
\]
The archive identifies this difference as a decisive snowball target for testing
whether structural noble-gas peaks and alkali-like drops survive promotion from
MDL proxy geometry to a physical Hamiltonian basis.
\end{definition}

\begin{obligation}[Structure-to-scale]
A complete physical promotion requires at least:
\begin{enumerate}
\item a defended Hamiltonian/operator carrier;
\item a scale/unit calibration;
\item a spectral observable;
\item valence/projector compatibility;
\item empirical comparison with declared error semantics.
\end{enumerate}
\end{obligation}

\section{Full recovery theorem shape}

\begin{definition}[Physical recovery package]
A full periodic-table recovery witness is a product of prior obligations:
\[
 \mathsf{PeriodicTableRecovery}
 =\mathsf{SpectrumRecovery}
 \times\mathsf{ShellRecovery}
 \times\mathsf{ValenceRecurrence}
 \times\mathsf{NuclearStabilityCompatibility}
 \times\mathsf{ObservableCompatibility}.
\]
Concretely it must supply:
\begin{enumerate}
\item a physical or formal spectrum for each admissible $Z$;
\item a shell decomposition with correct degeneracies and occupied projectors;
\item hard fermion occupation in $\bigwedge^Z\Hone$;
\item a ground-state/minimizer selector;
\item a robust projector/valence class;
\item recurrence across completion events;
\item nuclear-state compatibility for claims about real elements;
\item observable comparison under a declared calibration/error contract.
\end{enumerate}
\end{definition}

\begin{definition}[Generic element constructor]
A structural atomic generator has the type
\[
 \boxed{\Elem:\Nat_{>0}\longrightarrow\mathsf{AtomicState}}
\]
and, for neutral atoms, takes the schematic form
\[
 \Elem(Z)=
 (Z,N_Z^\star,Z,H_Z,\Psi_Z^\star,C_Z^\star,\Val_Z,Q_Z).
\]
There is no special constructor for hydrogen, carbon, iron, uranium, or element
$172$; they are evaluations of the same parameterized map.  Whether a given $Z$
has a physically stable nucleus is a separate theorem.
\end{definition}

\section{Cross-pollination map}

The formalism is intentionally reusable.  Its strongest justified cross-lane
connections are:
\begin{itemize}
\item \textbf{Triadic/369 algebra:} substrate and multiscale lift structure;
\item \textbf{observer/quotient machinery:} chemistry is a stable observable
quotient rather than the raw microscopic state;
\item \textbf{MDL/Lyapunov machinery:} selection and closure are phrased as
cost descent and resistance to admissible extension;
\item \textbf{representation theory:} shell degeneracy is generated by the
rotational/spin representation before filling;
\item \textbf{nuclear magic numbers:} same mathematical pattern---discrete
fermionic modes, exclusion, closure, enhanced stability---on a different
kernel; no identity of atomic and nuclear spectra is assumed;
\item \textbf{molecular assembly:} bonds become admissible shared-boundary
merges whose physical energy semantics remain separately payable;
\item \textbf{provenance hyperfabric:} every promotion remembers whether its
source was archive testimony, programmatic regression, exact formal theorem,
or external empirical evidence.
\end{itemize}

\section{Comparison with a finite topological matrix}

Suppose an external proposal provides a compact manifold $M$, a finite set of
geometric states $X(M)$, and a map
\[
 \chi:X(M)\to\{1,\ldots,Z_{\max}\}.
\]
Such a construction can be mathematically interesting, but it is a different
kind of object from the DASHI constructor unless it supplies a commuting bridge.

\begin{obligation}[External same-object bridge]
To identify the external chart with the DASHI atom generator, one needs maps
making at least the following structures agree:
\[
\begin{array}{ccc}
\T & \xrightarrow{\Phi} & \Rphys\\
\downarrow && \downarrow\\
X(M) & \xrightarrow{\chi} & \mathsf{AtomicState}
\end{array}
\]
together with preservation of fermionic admissibility, shell/valence recurrence,
and the physical observables used for validation.
\end{obligation}

In particular,
\[
 86+86=172
 \not\Rightarrow
 \text{a physically stable }Z=172\text{ nucleus}.
\]
The latter requires a nuclear-stability witness and downstream electronic
physical observables.

\section{Agda-shaped owner}

A compact Agda-facing owner can be organized as follows.  This is schematic
surface syntax rather than a claim that every field below is already inhabited.

\begin{verbatim}
record AtomicDictionary : Set1 where
  field
    protonNumber   : Nat
    neutronNumber  : Nat
    electronNumber : Nat

    Kernel         : Set
    Mode           : Kernel -> Set
    occupied       : Kernel -> Bool
    pauliOK        : Set

    Boundary       : Set
    active         : Boundary -> Set

    Cost           : Set
    lessCost       : Cost -> Cost -> Set
    mdl            : Kernel -> Cost

    ValenceClass   : Set
    valence        : ValenceClass

record PeriodicTableConstructor : Set1 where
  field
    atom              : Nat -> AtomicDictionary
    completionStep    : Nat -> Nat -> Set
    recurrentValence  : completionStep -> Set
    periodicQuotient  : Nat -> Nat -> Set
    nonClaimBoundary  : List String
\end{verbatim}

\section{Non-collapse laws}

The theory is defined as much by what it refuses as by what it constructs:
\[
\boxed{
\begin{aligned}
369\text{ cardinality} &\not\Rightarrow \text{orbital quantum numbers},\\
\text{shell capacity} &\not\Rightarrow \text{energy ordering},\\
\text{MDL ordering} &\not\Rightarrow \text{physical Hamiltonian spectrum},\\
\text{electronic filling} &\not\Rightarrow \text{nuclear stability},\\
Z &\not\Rightarrow \text{stable isotope},\\
\text{closed shell} &\not\Rightarrow \text{quantitative ionization energy},\\
\text{valence recurrence} &\not\Rightarrow \text{finished chemistry},\\
\text{geometric chart} &\not\Rightarrow \text{same-object atomic constructor}.
\end{aligned}}
\]

\section{Final normal form}

The normal form is the commuting dependency diagram:
\[
\boxed{
\begin{array}{c}
\text{369 / ternary carrier}\\
\downarrow\\
\text{contraction / stable structural classes}\\
\downarrow\\
\text{physical realization}\\
\downarrow\\
\text{one-particle kernel representation}\\
\downarrow\\
\text{spin + angular degeneracy}\\
\downarrow\\
\text{fermionic exterior power / exclusion quotient}\\
\downarrow\\
\text{MDL or Hamiltonian minimizer}\\
\downarrow\\
\text{shell occupancy}\\
\downarrow\\
\text{energetic/writable boundary projection}\\
\downarrow\\
\text{valence recurrence across shell completion}\\
\downarrow\\
\text{periodic-table quotient}\\
\downarrow\\
\text{chemical observables and validation}.
\end{array}}
\]

Equivalently:
\[
\boxed{\begin{gathered}
\text{geometry/algebra generates the admissible state space;}\\
\text{fermionic MDL dynamics selects the atom;}\\
\text{valence recurrence generates the periodic table.}
\end{gathered}}
\]

\section{Chronology, publication identity, and verification status}

The chronology below is a provenance surface, not a premise of the mathematical
construction.  Dates name the earliest clean receipts located in this audit;
they do not retroactively date conception, and a repository commit is not treated
as an external publication.

\begin{longtable}{p{0.16\textwidth}p{0.26\textwidth}p{0.46\textwidth}}
\toprule
Date & Surface & Status / significance\\
\midrule
2025-11-11 & \texttt{dashifine} spectral-line tooling & Earliest clean atomic/spectral implementation located in the wider repo archaeology: commit \texttt{4f1441e4989b...}, ``Add spectral line tooling and smoke tests.'' This predates the direct DASHI atomic constructor but is not itself the 369 periodic-table formalism.\\
2026-02-25 / 2026-03-05 & \texttt{dashifine} geometry validation & Authored 25 February and committed 5 March: topology along the RG/projection axis plus low-dimensional-manifold diagnostics. This is a programme precursor, not atom recovery.\\
2026-03-06 & \texttt{dashiQ/PHYSICS\_TARGETS.md} & Explicit projection/effective-manifold programme, including emergent $3+1$D/Lorentz/locality questions. Again a programme precursor rather than direct atomic construction.\\
2026-04-30 21:51:43 +10:00 & First direct atom/chemistry repo implementation located & \texttt{dashi\_agda} commit \texttt{42e1d740141b0e9e1ca717ae5df79d6e31546c07}: atom/chemistry recovery carrier plus closed-shell stability and shell-filling dictionary compatibility. The commit itself says this is staged closure, not finished chemistry recovery.\\
2026-05-03 & Explicit atomic generation pipeline & Follow-up commits add an explicit atomic generation pipeline and thread chemistry into it. This is the direct generative-constructor strengthening of the April carrier.\\
2026-08-12 01:33:28Z & Lean 369 mirror at its current path & \texttt{chboishabba/dashi\_lean4} commit \texttt{72734285fd83387837e0025eb51a93b63629a0b9} reorganizes the Aristotle artifacts under \texttt{AgdaMirror/}; \texttt{AgdaMirror/Base369.lean} explicitly contains proofs of ternary/hexadic/nonary rotation/XOR laws. This is substrate-level cross-assistant proof, not an atom constructor mirror.\\
2026-09-11 & Current formal consolidation & Draft PR \#886 adds the generative 369 atom owner, provenance/snowball owner, chronology/status owner, focused validation root, and this manuscript.\\
\bottomrule
\end{longtable}

The attached historical conversation export visibly prints the page label
``12/1/26''.  That is retained as a source-object date cue only: the locale and
original publication semantics are not independently established, so this note
does not silently convert it into an asserted ISO publication date.

\subsection{First implementation versus first publication}

The strongest date claim paid by the current repository archaeology is therefore:
\[
\boxed{\text{first direct repo implementation located} = \text{30 April 2026}.}
\]
The manuscript does \emph{not} currently have a same-object DOI, arXiv, Zenodo,
or other external-publication receipt predating this PR.  Accordingly:
\[
\boxed{\text{repo implementation date}\neq\text{external publication date}.}
\]
The present paper is the first repo-contained manuscript located in this audit
that is devoted specifically to the reconstructed generative atom/periodic-table
formalism.  That is a manuscript/provenance claim, not a priority claim outside
the repositories.

\subsection{What has actually been checked or executed}

The project distinguishes mathematical status from certification status:

\begin{longtable}{p{0.25\textwidth}p{0.22\textwidth}p{0.43\textwidth}}
\toprule
Object & Current status & Evidence and boundary\\
\midrule
Historical MDL/exclusion filling regression & Reported executed run; artifact missing & The attached archive reports a programmatic run over the first 15+ elements and closure coordinates $Z=2,10,18$.  The exact original Python file, parameter schedule, commit/hash, and rerun receipt have not yet been located, so the manuscript does not call this reproduced.\\
Existing Agda atom/369 owners on \texttt{master} & Source/theorem terms present; not freshly checked in this session & The repo has an explicit validation policy: a targeted pass is an Agda compile signal and aggregate green means \texttt{agda -i . DASHI/Everything.agda} passes.  The execution environment used for this paper has no Agda binary, so no fresh local compiler receipt is claimed.\\
PR \#886 Agda modules & Authored, validation-rooted, pending compiler receipt & The PR contains the generative owner, provenance owner, chronology/status owner, and focused validation root.  No PR workflow run had appeared at the inspected head, and no local Agda executable was available.  Source presence is not typechecking.\\
\texttt{dashi\_lean4/AgdaMirror/Base369.lean} & Explicit proof terms in source; independent build receipt not obtained here & The file describes itself as a genuine fully proved mirror and contains concrete proofs by case analysis for the $3/6/9$ carrier laws.  The repository contains a \texttt{leanprover/lean-action} CI workflow, but the inspected main head returned no workflow runs/status checks, and this session did not independently rebuild Mathlib/Lean.\\
Full empirical periodic-table recovery & Formal interface / open obligation & The current theory separates spectra, nuclear stability, scale calibration, ionization observables, bonding, and empirical validation.  Their typed interfaces must not be confused with physically inhabited recovery theorems.\\
This \LaTeX{} manuscript & Compiled and rendered & The source is compiled to the accompanying PDF and rendered for layout inspection.  Successful document rendering is only a publication-artifact check, not a proof-assistant or empirical-physics check.\\
\bottomrule
\end{longtable}

Thus the public status is intentionally multi-layered:
\[
\boxed{
\text{mathematical specification}
\neq \text{proof-assistant source}
\neq \text{recorded typecheck/build}
\neq \text{executed regression}
\neq \text{empirical physical recovery}.}
\]

\subsection{Lean cross-pollination}

The Lean repository adds a useful independent substrate lane without licensing an
overclaim.  Its \texttt{AgdaMirror/Base369.lean} defines \texttt{TriTruth},
\texttt{HexTruth}, and \texttt{NonaryTruth}; defines the iterated rotation
\texttt{spin}; and proves the agreement of spin-defined and modular XOR, the
respective rotation orders, identities, and ternary associativity.  This means
that the elementary 369 algebra is not Agda-only source prose.  What was not found
in the current Lean archaeology is a direct Lean mirror of the atomic shell,
valence, or periodic-table constructor itself.  The safe cross-language statement
is therefore
\[
\boxed{\text{Lean-paid 369 substrate}\quad\text{but}\quad
\text{atomic constructor presently Agda-facing}.}
\]

\appendix
\section{Provenance notes from the attached archive}

The attached \emph{DASHI Atom} archive records the following source claims used
above.

\begin{enumerate}
\item Noble gases are treated as saturated-kernel objects with maximal symmetry,
minimal interaction, and MDL resistance to extension.
\item Pages 294--296 state the remembered programmatic path: kernel states
$(n,\ell,m,s)$, exclusion quotient, MDL proxy $E(n,\ell)=n+\alpha\ell$, closures
at $Z=2,10,18$, and structure/scale separation.
\item Pages 297--301 provide the atom-centric reconstruction: atoms are stable
filling configurations of a discrete fermionic kernel; exclusion forces finite
occupancy; MDL ordering selects fillings; noble gases are local MDL minima;
groups are boundary-condition equivalence classes; nuclear magic numbers share
the same mathematics at another kernel.
\item Page 418 records the scale snowball: ionization energy
$IE_1(Z)=E(Z-1)-E(Z)$ is the next validation route for converting MDL geometry
into calibrated physical observables.
\end{enumerate}

\end{document}
